
\begin{frame}[fragile]{Interrogation}
  \begin{redbullet}
    \item ElasticSearch s’appuie sur le système d’indexation \textbf{Lucene}

\item Lucene propose un langage de recherche basé sur des combinaisons de
  mot-clés, langage étendu et raffiné par ElasticSearch (cf plus tard)
\item La première méthode pour transmettre des recherches est de passer une
  expression en paramètre à l’URL :
\begin{minted}[frame=lines,
               baselinestretch=.9,
               framesep=2mm]{bash}
https://localhost:9200/films/_search?q=Rome
http://localhost:9200/films/_search?q=Rome,Paris
http://localhost:9200/films/_search?q=1994
# utilisable dans l'interface ElasticVue, onglet Rest
\end{minted}

\end{redbullet}
\end{frame}

\begin{frame}[fragile]{la réponse d'ES}
	\begin{minted}[ baselinestretch=0.9, fontsize=\footnotesize, ]{bash}
{
  "took" : 1,
  "timed_out" : false,
  "_shards" : {
    "total" : 1,
    "successful" : 1,
    "failed" : 0
  },
  "hits" : {
    "total" : 1,
    "max_score" : 0.558217,
    "hits" : [ {
      "_index" : "nfe204",
      "_type" : "movies",
      "_id" : "movie:2",
      "_score" : 0.558217,
      "_source":{ "title": "romz, città aperta", "year": 1946, 
                  "genre": "Drama",
                   "summary": "... ", 
                   "country": "Italy", 
                   "director": "R. Rosselini", 
				   }
    } ]
  }
}

\end{minted} 
\end{frame} 


\begin{frame}[fragile]{Termes}

  \begin{redbullet}
    \item Notion de base : le {\bf terme}
    \item c'est un mot au sens usuel
    \item ou une séquence de mots entre apostrophes
  \end{redbullet}
\smallskip
On peut interroger un index avec :
\begin{minted}[frame=lines, baselinestretch=.9, fontsize=\footnotesize, framesep=2mm]{sql}
Neo apprend
\end{minted}
Puis :
\begin{minted}[frame=lines, baselinestretch=.9, fontsize=\footnotesize, framesep=2mm]{sql}
  "Neo apprend"
\end{minted}
\vfill
\begin{blocgauche}
  Première recherche : documents avec ``Neo'', ``apprend'' ou les deux
\end{blocgauche}

\end{frame}

\begin{frame}[fragile]{Termes (suite)}
  \begin{redbullet}
    \item la recherche d'un terme s'effectue toujours sur un champ. 
\item La syntaxe complète pour associer le champ et le terme est:
\begin{minted}[frame=lines,
               baselinestretch=.9,
               fontsize=\footnotesize,
               framesep=2mm]{sql}
  champ:terme
\end{minted}
\item si non précisé, c'est le champ par défaut qui est utilisé
\item pratique courante : concaténer toutes les chaînes de caractères en un
  champ ``\_all'' général, défini par défaut
\item Nos requêtes deviennent :
\begin{minted}[baselinestretch=.9,
               fontsize=\footnotesize,
               framesep=2mm]{sql}

               _all:Neo _all:apprend
\end{minted}
\item et 
\begin{minted}[baselinestretch=.9,
               fontsize=\footnotesize,
               framesep=2mm]{sql}
               _all:"Neo apprend"
\end{minted}
  \end{redbullet}
\end{frame}

\begin{frame}[fragile]{Termes (suite)}

  \begin{redbullet}
    \item Les valeurs des termes (dans la requête) et le texte indexé sont tous deux
soumis à des transformations spécifiées dans le schéma. 
\item Une transformation simple est de tout transcrire en minuscules. 
\begin{minted}[frame=lines,
               baselinestretch=.9,
               fontsize=\footnotesize,
               framesep=2mm]{sql}
              _all:"Neo APPREND"
\end{minted}
\item  Les transformations appliquées à la requête ET au texte indexé doivent être 
   cohérentes : si les termes sont transformés en majuscules, et le texte indexé en minuscules,
   on n'aura jamais de résultat!
  \end{redbullet}
\end{frame}

\begin{frame}[fragile]{Termes (suite)}
  On peut spécifier des termes (pas des séquences) incomplets
  \begin{redbullet}
      \setlength{\itemsep}{8pt}
    \item le '?' indique un caractère inconnu;  ``opti?al'' désigne ``optimal'', ``optical'', etc.
      
    \item le '*' indique n'importe quelle séquence de caractères; 
      ``opti*'' pour toute chaîne commençant par ``opti''
    \item Approximation avec $\sim$: rechercher ``optimal`` et ``optimal$\sim$`` (0 et 1 résultat, ''optical'')
    
       La similarité est évaluée par une distance d'édition 
      (nb opérations pour passer de ``optimal'' à ``optical'')
  
    \item Les intervalles sont exprimés avec  $[ ]$ (bornes comprises) ou  \{ \} bornes exclues
  \end{redbullet}
\begin{minted}[frame=lines,
               baselinestretch=.9,
               fontsize=\footnotesize,
               framesep=2mm]{sql}
    year:[1984 TO 2010]
\end{minted}
\end{frame}

\begin{frame}[fragile]{Requêtes Booléennes}

  \begin{redbullet}
    \item Les critères peuvent être combinés avec des \alert{opérateurs
      Booléens} :\\
      \variable{AND}, \variable{OR} et \variable{NOT}
    \item Attention : majuscules
  \end{redbullet}
\begin{minted}[frame=lines,
               baselinestretch=.9,
               fontsize=\footnotesize,
               framesep=2mm]{sql}
year:[1990 TO 2005] OR title:M*
year:[1990 TO 2005] AND NOT title:M*
\end{minted}
   
\begin{redbullet}
  \item Par défaut, c'est un \variable{OR} qui est appliqué
  \item Recherche sur plusieurs critères ramène l'union des résultats sur chaque
    critère pris individuellement
  %\item La requête suivante recherche les produits Dell {\bf ou} dont la popularité est
%égale à 6 :
%\begin{minted}[frame=lines,
               %baselinestretch=.9,
               %fontsize=\footnotesize,
               %framesep=2mm]{sql}
    %%popularity:6  manu:Dell
%\end{minted}
\end{redbullet}
\end{frame}

\begin{frame}[fragile]{Requêtes structurées}

On fournit un document décrivant la requête en la \textit{postant}  
(méthode \texttt{POST}):

\vfill

\begin{minted}[ baselinestretch=0.9]{bash}
	{
    	"query": {
      		"match": {
        		"title": "Star Wars"
      		}
  		}
   	}
\end{minted}

\vfill

Un premier pas vers une interrogation plus structurée.

\end{frame}

\begin{frame}[fragile]{Choix des champs et du résultat}

On peut ajouter des spécifications dans la requête.

\vfill

\begin{minted}[ baselinestretch=0.9]{bash}
	{
    	"query": {
      		"match": {
        		"title": "Star Wars"
      		}
  		},
  		"fields": ["title"],
  		"_source": false
   	}
\end{minted}

\end{frame}

\begin{frame}[fragile]{Effectuer des recherches approchées}

\begin{minted}[ baselinestretch=0.9]{bash}
	{
	  "query": {
	     "fuzzy": {
      		"title": {
        		"value": "matrice"
    	  	}
    	  }
  		},
  	 "fields": ["title"],
	 "_source": false
	}
\end{minted}

\end{frame}


\begin{frame}[fragile]{Recherches par intervalle}

\begin{minted}[ baselinestretch=0.9]{bash}
	{
  		"query": {
		    "range": {
		      "year": {"gte": 2010, "lte": 2020}
    		}
  		},
  		"fields": ["title", "year"],
  		"_source": false
	}
\end{minted}

\end{frame}





%\begin{frame}[fragile]{Opérateur $+$, classement}
  %\begin{redbullet}
    %\item préfixe d'un nom de champ, il indique que la valeur du champ {\bf doit}
      %être égale au terme
    %\item il existe également un opérateur $-$, équivalent au \variable{NOT}
    %\item recherche des documents dont la popularité est 6 (obligatoire) et qui peuvent
%être produits par Dell ou un autre constructeur 
%\begin{minted}[frame=lines,
               %baselinestretch=.9,
               %fontsize=\footnotesize,
               %framesep=2mm]{sql}
    %%+popularity:6  manu:Dell
%\end{minted}
%\item différence avec ce qui précède : le {\bf classement} du moteur
%\item illustre la différence entre recherche d'information et interrogation
  %de bases de données
  %\end{redbullet}
  %\bigskip
  %Interpréter un classement est parfois délicat :\\
%\begin{minted}[frame=lines,
               %baselinestretch=.9,
               %fontsize=\footnotesize,
               %framesep=2mm]{sql}
    %%+popularity:6  cat:electronics

    %%+popularity:6  -manu:Dell

%\end{minted}
%\end{frame}


\begin{frame}[fragile]{En conclusion}


Ce qui précère constitue l'essentiel des opérations d'interrogation avec un moteur
de recherche.

  \begin{redbullet}
    \item On est loin de la finesse d'interrogation d'un langage comme SQL.
    
    \item On fait peu mais on fait bien: systèmes extrêmement efficaces, même pour
        de très grands volumes. 
  \end{redbullet}

\vfill

\alert{Surtout}, deux différences essentielles:

 \begin{redbullet}
    \item Le système permet des recherches \alert{approchées}: on trouve tous les documents "proches de" la requête
    \item Le résultat peut être \alert{ordonné} par pertinence décroissante.
  \end{redbullet}

\vfill

Comme ça marche? Voir la suite...
\end{frame}

